\exnumber{04}
\extitle{UX IV - Ex 5 - Prototypage}
\exfiles{ex04/fichier Figma, nomme Name\_FirstName\_DeliverableName\_Date\_VersionNumber}
\exauthorize{Ceux nommés dans l'exercice}
\makeheaderfiles

\textbf{Instructions :}

Donnez une autre dimension à votre user flow et transformez votre user flow en prototype papier (12 écrans max).

D'abord, dessinez des prototypes basse fidélité sur papier. Un écran principal = 1 frame mobile.

Ensuite, dessinez vos écrans mid-fidelity sur Figma avec des formes simples, en noir et blanc.

\textbf{Pistes :} Vocabulaire

\begin{itemize}
\item Basse fidélité : prototypes papier - exemple \url{https://miro.medium.com/max/1400/1*5nUfqqA2gjdAYHagjbHA5w.jpeg}
\item Moyenne fidélité : wireframes, noir et blanc, seulement les éléments principaux, sur un outil de prototype - exemple \url{https://miro.medium.com/max/1400/1*eGSfaxwYbxUFomYbyO6GWw.png}
\item Haute fidélité : wireframes avec couleur, styles, détails graphiques et micro-interactions - exemple \url{https://miro.medium.com/max/1400/1*Xn0HSKAvhr4TZzC9lN5udw.gif}
\end{itemize}

\textbf{Livrables :}

\begin{itemize}
\item Photo d'un prototype papier
\item Un fichier Figma avec les wireframes (moyenne fidélité)
\item Déposez les documents sur Github, dans un zip nommé « ProjectTitle\_Name\_FirstName », et chaque livrable ainsi : « Name\_FirstName\_DeliverableName\_Date\_VersionNumber ». Exemple : « Doe\_John\_MidFiPrototype\_05242024\_V1 ».
\item Ajoutez un titre dans tous les documents écrits.
\item Soignez l'aspect global des livrables. Ils doivent être clairs, simples et lisibles. Vous pouvez regarder canva.com, sans surcharger le dessin de détails.
\end{itemize}

\textbf{Ressources :}

\begin{itemize}
\item Exemple : Prototypage papier - Le guide pratique pour débutants \url{https://www.uxpin.com/studio/blog/paper-prototyping-the-practical-beginners-guide/}
\item Exemple : Application Wellness \url{https://medium.com/@ebtaskin/story-of-my-very-first-ux-ui-design-project-at-ironhack-bootcamp-955addeb760e}
\item Template Mobile Android \url{https://i.pinimg.com/originals/43/8e/de/438ede24264eb9552b635d34142eac29.png}
\item Template iPhone \url{https://www.cssauthor.com/wp-content/uploads/2014/01/iPhone-5S-Template-PSD.jpg}
\item Figma
\end{itemize}

\textbf{Citations :}

\begin{itemize}
\item « Si vous n'êtes pas gêné par la première version de votre produit, vous l'avez lancé trop tard » Reid Hoffman | LinkedIn
\item Pour prototyper votre solution, vous aurez besoin d'un changement temporaire de philosophie : du parfait au juste suffisant, de la qualité à long terme à la simulation temporaire.
\item Le prototype sert à répondre à des questions, donc restez concentré. Vous avez juste besoin d'une façade réaliste à laquelle les clients peuvent réagir.
\end{itemize}

